\textbf{REFERENCES}

\href{https://sciwheel.com/work/bibliography/1306567}{1.~   Chang, J.T.,
Wherry, E.J., and Goldrath, A.W. (2014). Molecular regulation of
effector and memory T cell differentiation. Nat. Immunol. \emph{15},
1104--1115. 10.1038/ni.3031.}

\href{https://sciwheel.com/work/bibliography/15592036}{2.~   Koh, C.-H.,
Lee, S., Kwak, M., Kim, B.-S., and Chung, Y. (2023). CD8 T-cell subsets:
heterogeneity, functions, and therapeutic potential. Exp. Mol. Med.
\emph{55}, 2287--2299. 10.1038/s12276-023-01105-x.}

\href{https://sciwheel.com/work/bibliography/8929852}{3.~   Adams, N.M.,
Grassmann, S., and Sun, J.C. (2020). Clonal expansion of innate and
adaptive lymphocytes. Nat. Rev. Immunol. \emph{20}, 694--707.
10.1038/s41577-020-0307-4.}

\href{https://sciwheel.com/work/bibliography/100389}{4.~   Masopust, D.,
and Schenkel, J.M. (2013). The integration of T cell migration,
differentiation and function. Nat. Rev. Immunol. \emph{13}, 309--320.
10.1038/nri3442.}

\href{https://sciwheel.com/work/bibliography/18537397}{5.~   Lavenant,
H., Zhang, S., Kim, Y.-H., and Schiebinger, G. (2021). Towards a
mathematical theory of trajectory inference. arXiv.
10.48550/arxiv.2102.09204.}

\href{https://sciwheel.com/work/bibliography/10552252}{6.~   Huang, H.,
Zhou, P., Wei, J., Long, L., Shi, H., Dhungana, Y., Chapman, N.M., Fu,
G., Saravia, J., Raynor, J.L., et al. (2021). In~vivo CRISPR screening
reveals nutrient signaling processes underpinning CD8+ T~cell fate
decisions. Cell \emph{184}, 1245-1261.e21. 10.1016/j.cell.2021.02.021.}

\href{https://sciwheel.com/work/bibliography/7598004}{7.~   Yoon, H.,
Kim, T.S., and Braciale, T.J. (2010). The cell cycle time of CD8+ T
cells responding in vivo is controlled by the type of antigenic
stimulus. PLoS ONE \emph{5}, e15423. 10.1371/journal.pone.0015423.}

\href{https://sciwheel.com/work/bibliography/18384638}{8.~   He, S.,
Zhu, Y., Tavakol, D.N., Ye, H., Lao, Y.-H., Zhu, Z., Xu, C., Chauhan,
S., Garty, G., Tomer, R., et al. (2026). Squidiff: predicting cellular
development and responses to perturbations using a diffusion model. Nat.
Methods \emph{23}, 65--77. 10.1038/s41592-025-02877-y.}

\href{https://sciwheel.com/work/bibliography/17494421}{9.~   Park, C.,
Mani, S., Beltran-Velez, N., Maurer, K., Huang, T., Li, S., Gohil, S.,
Livak, K.J., Knowles, D.A., Wu, C.J., et al. (2024). A Bayesian
framework for inferring dynamic intercellular interactions from
time-series single-cell data. Genome Res. \emph{34}, 1384--1396.
10.1101/gr.279126.124.}

\href{https://sciwheel.com/work/bibliography/18201304}{10.~  Bailey,
S.R., Bartee, E., Daniels, K.G., Heery, C.R., Kaumaya, P., Lesinski,
G.B., Lowinger, T.B., Nelson, M.H., Rubinstein, M.P., Wittling, M.C., et
al. (2025). Constructing the cure: engineering the next wave of antibody
and cellular immune therapies. J. Immunother. Cancer \emph{13}.
10.1136/jitc-2025-011761.}

\href{https://sciwheel.com/work/bibliography/17904669}{11.~  Green,
W.D., Gomez, A., Plotkin, A.L., Pratt, B.M., Merritt, E.F., Mullins,
G.N., Kren, N.P., Modliszewski, J.L., Zhabotynsky, V., Woodcock, M.G.,
et al. (2025). Enhancer-driven gene regulatory networks reveal
transcription factors governing T cell adaptation and differentiation in
the tumor microenvironment. Immunity \emph{58}, 1725-1741.e9.
10.1016/j.immuni.2025.04.030.}

\href{https://sciwheel.com/work/bibliography/8940021}{12.~  Milner,
J.J., Toma, C., He, Z., Kurd, N.S., Nguyen, Q.P., McDonald, B., Quezada,
L., Widjaja, C.E., Witherden, D.A., Crowl, J.T., et al. (2020).
Heterogenous Populations of Tissue-Resident CD8+ T Cells Are Generated
in Response to Infection and Malignancy. Immunity \emph{52}, 808-824.e7.
10.1016/j.immuni.2020.04.007.}

\href{https://sciwheel.com/work/bibliography/8917053}{13.~  Kurd, N.S.,
He, Z., Louis, T.L., Milner, J.J., Omilusik, K.D., Jin, W., Tsai, M.S.,
Widjaja, C.E., Kanbar, J.N., Olvera, J.G., et al. (2020). Early
precursors and molecular determinants of tissue-resident memory CD8+ T
lymphocytes revealed by single-cell RNA sequencing. Sci. Immunol.
\emph{5}. 10.1126/sciimmunol.aaz6894.}

\href{https://sciwheel.com/work/bibliography/15807051}{14.~  Chung,
H.K., Liu, C., Battu, A., Jambor, A.N., Pratt, B.M., Xie, F.,
Riesenberg, B.P., Casillas, E., Sun, M., Landoni, E., et al. (2025).
Atlas-Guided Discovery of Transcription Factors for T Cell Programming.
BioRxiv. 10.1101/2023.01.03.522354.}

\href{https://sciwheel.com/work/bibliography/14885142}{15.~  Chu, Y.,
Dai, E., Li, Y., Han, G., Pei, G., Ingram, D.R., Thakkar, K., Qin,
J.-J., Dang, M., Le, X., et al. (2023). Pan-cancer T cell atlas links a
cellular stress response state to immunotherapy resistance. Nat. Med.
\emph{29}, 1550--1562. 10.1038/s41591-023-02371-y.}

\href{https://sciwheel.com/work/bibliography/16689366}{16.~  Buquicchio,
F.A., Fonseca, R., Yan, P.K., Wang, F., Evrard, M., Obers, A.,
Gutierrez, J.C., Raposo, C.J., Belk, J.A., Daniel, B., et al. (2024).
Distinct epigenomic landscapes underlie tissue-specific memory T~cell
differentiation. Immunity \emph{57}, 2202-2215.e6.
10.1016/j.immuni.2024.06.014.}

\href{https://sciwheel.com/work/bibliography/5496398}{17.~  Street, K.,
Risso, D., Fletcher, R.B., Das, D., Ngai, J., Yosef, N., Purdom, E., and
Dudoit, S. (2018). Slingshot: cell lineage and pseudotime inference for
single-cell transcriptomics. BMC Genomics \emph{19}, 477.
10.1186/s12864-018-4772-0.}

\href{https://sciwheel.com/work/bibliography/2919207}{18.~  Haghverdi,
L., Büttner, M., Wolf, F.A., Buettner, F., and Theis, F.J. (2016).
Diffusion pseudotime robustly reconstructs lineage branching. Nat.
Methods \emph{13}, 845--848. 10.1038/nmeth.3971.}

\href{https://sciwheel.com/work/bibliography/9367004}{19.~  Bergen, V.,
Lange, M., Peidli, S., Wolf, F.A., and Theis, F.J. (2020). Generalizing
RNA velocity to transient cell states through dynamical modeling. Nat.
Biotechnol. \emph{38}, 1408--1414. 10.1038/s41587-020-0591-3.}

\href{https://sciwheel.com/work/bibliography/13729536}{20.~  Hudson,
W.H., and Wieland, A. (2023). Technology meets TILs: Deciphering T~cell
function in the -omics era. Cancer Cell \emph{41}, 41--57.
10.1016/j.ccell.2022.09.011.}

\href{https://sciwheel.com/work/bibliography/4615200}{21.~  Youngblood,
B., Hale, J.S., Kissick, H.T., Ahn, E., Xu, X., Wieland, A., Araki, K.,
West, E.E., Ghoneim, H.E., Fan, Y., et al. (2017). Effector CD8 T cells
dedifferentiate into long-lived memory cells. Nature \emph{552},
404--409. 10.1038/nature25144.}

\href{https://sciwheel.com/work/bibliography/5308353}{22.~  Herndler-Brandstetter,
D., Ishigame, H., Shinnakasu, R., Plajer, V., Stecher, C., Zhao, J.,
Lietzenmayer, M., Kroehling, L., Takumi, A., Kometani, K., et al.
(2018). KLRG1+ Effector CD8+ T Cells Lose KLRG1, Differentiate into All
Memory T Cell Lineages, and Convey Enhanced Protective Immunity.
Immunity \emph{48}, 716-729.e8. 10.1016/j.immuni.2018.03.015.}

\href{https://sciwheel.com/work/bibliography/637212}{23.~  Masopust, D.,
Choo, D., Vezys, V., Wherry, E.J., Duraiswamy, J., Akondy, R., Wang, J.,
Casey, K.A., Barber, D.L., Kawamura, K.S., et al. (2010). Dynamic T cell
migration program provides resident memory within intestinal epithelium.
J. Exp. Med. \emph{207}, 553--564. 10.1084/jem.20090858.}

\href{https://sciwheel.com/work/bibliography/6347423}{24.~  Schiebinger,
G., Shu, J., Tabaka, M., Cleary, B., Subramanian, V., Solomon, A.,
Gould, J., Liu, S., Lin, S., Berube, P., et al. (2019).
Optimal-Transport Analysis of Single-Cell Gene Expression Identifies
Developmental Trajectories in Reprogramming. Cell \emph{176},
928-943.e22. 10.1016/j.cell.2019.01.006.}

\href{https://sciwheel.com/work/bibliography/15548348}{25.~  Bunne, C.,
Meng-Papaxanthos, L., Krause, A., and Cuturi, M. (2021). Proximal
Optimal Transport Modeling of Population Dynamics. arXiv.
10.48550/arxiv.2106.06345.}

\href{https://sciwheel.com/work/bibliography/17428135}{26.~  Klein, D.,
Palla, G., Lange, M., Klein, M., Piran, Z., Gander, M.,
Meng-Papaxanthos, L., Sterr, M., Saber, L., Jing, C., et al. (2025).
Mapping cells through time and space with moscot. Nature \emph{638},
1065--1075. 10.1038/s41586-024-08453-2.}

\href{https://sciwheel.com/work/bibliography/14086393}{27.~  Gorin, G.,
Fang, M., Chari, T., and Pachter, L. (2022). RNA velocity unraveled.
PLoS Comput. Biol. \emph{18}, e1010492. 10.1371/journal.pcbi.1010492.}

\href{https://sciwheel.com/work/bibliography/17425314}{28.~  Reina-Campos,
M., Monell, A., Ferry, A., Luna, V., Cheung, K.P., Galletti, G.,
Scharping, N.E., Takehara, K.K., Quon, S., Challita, P.P., et al.
(2025). Tissue-resident memory CD8 T cell diversity is spatiotemporally
imprinted. Nature \emph{639}, 483--492. 10.1038/s41586-024-08466-x.}

\href{https://sciwheel.com/work/bibliography/14778778}{29.~  Gill, A.L.,
Hudson, W.H., Valanparambil, R.M., Ahn, E., McGuire, D.J., Wieland, A.,
McManus, D.T., Kissick, H.T., Akondy, R.S., and Ahmed, R. (2023).
Longitudinal Analysis of the Phenotype, Transcriptional Profile, and
Anatomic Location of Memory CD8 T Cell Subsets after Acute Viral
Infection. J. Virol. \emph{97}, e0155622. 10.1128/jvi.01556-22.}

\href{https://sciwheel.com/work/bibliography/1124271}{30.~  Zhang, N.,
and Bevan, M.J. (2011). CD8(+) T cells: foot soldiers of the immune
system. Immunity \emph{35}, 161--168. 10.1016/j.immuni.2011.07.010.}

\href{https://sciwheel.com/work/bibliography/13830433}{31.~  Giles,
J.R., Ngiow, S.F., Manne, S., Baxter, A.E., Khan, O., Wang, P., Staupe,
R., Abdel-Hakeem, M.S., Huang, H., Mathew, D., et al. (2022). Shared and
distinct biological circuits in effector, memory and exhausted CD8+ T
cells revealed by temporal single-cell transcriptomics and epigenetics.
Nat. Immunol. \emph{23}, 1600--1613. 10.1038/s41590-022-01338-4.}

\href{https://sciwheel.com/work/bibliography/808308}{32.~  Kaech, S.M.,
and Cui, W. (2012). Transcriptional control of effector and memory CD8+
T cell differentiation. Nat. Rev. Immunol. \emph{12}, 749--761.
10.1038/nri3307.}

\href{https://sciwheel.com/work/bibliography/425675}{33.~  Gerlach, C.,
Rohr, J.C., Perié, L., van Rooij, N., van Heijst, J.W.J., Velds, A.,
Urbanus, J., Naik, S.H., Jacobs, H., Beltman, J.B., et al. (2013).
Heterogeneous differentiation patterns of individual CD8+ T cells.
Science \emph{340}, 635--639. 10.1126/science.1235487.}

\href{https://sciwheel.com/work/bibliography/11099644}{34.~  Mold, J.E.,
Modolo, L., Hård, J., Zamboni, M., Larsson, A.J.M., Stenudd, M.,
Eriksson, C.-J., Durif, G., Ståhl, P.L., Borgström, E., et al. (2021).
Divergent clonal differentiation trajectories establish CD8+ memory
T~cell heterogeneity during acute viral infections in humans. Cell Rep.
\emph{35}, 109174. 10.1016/j.celrep.2021.109174.}

\href{https://sciwheel.com/work/bibliography/16029036}{35.~  Mold, J.E.,
Weissman, M.H., Ratz, M., Hagemann-Jensen, M., Hård, J., Eriksson,
C.-J., Toosi, H., Berghenstråhle, J., Ziegenhain, C., von Berlin, L., et
al. (2024). Clonally heritable gene expression imparts a layer of
diversity within cell types. Cell Syst. \emph{15}, 149-165.e10.
10.1016/j.cels.2024.01.004.}

\href{https://sciwheel.com/work/bibliography/1306566}{36.~  Roychoudhuri,
R., Lefebvre, F., Honda, M., Pan, L., Ji, Y., Klebanoff, C.A., Nichols,
C.N., Fourati, S., Hegazy, A.N., Goulet, J.-P., et al. (2015).
Transcriptional profiles reveal a stepwise developmental program of
memory CD8(+) T cell differentiation. Vaccine \emph{33}, 914--923.
10.1016/j.vaccine.2014.10.007.}

\href{https://sciwheel.com/work/bibliography/13843025}{37.~  Daniel, B.,
Yost, K.E., Hsiung, S., Sandor, K., Xia, Y., Qi, Y., Hiam-Galvez, K.J.,
Black, M., J Raposo, C., Shi, Q., et al. (2022). Divergent clonal
differentiation trajectories of T cell exhaustion. Nat. Immunol.
\emph{23}, 1614--1627. 10.1038/s41590-022-01337-5.}

\href{https://sciwheel.com/work/bibliography/14304004}{38.~  Quezada,
L.K., Jin, W., Liu, Y.C., Kim, E.S., He, Z., Indralingam, C.S., Tysl,
T., Labarta-Bajo, L., Wehrens, E.J., Jo, Y., et al. (2023). Early
transcriptional and epigenetic divergence of CD8+ T cells responding to
acute versus chronic infection. PLoS Biol. \emph{21}, e3001983.
10.1371/journal.pbio.3001983.}

\href{https://sciwheel.com/work/bibliography/12129174}{39.~  Gearty,
S.V., Dündar, F., Zumbo, P., Espinosa-Carrasco, G., Shakiba, M.,
Sanchez-Rivera, F.J., Socci, N.D., Trivedi, P., Lowe, S.W., Lauer, P.,
et al. (2022). An autoimmune stem-like CD8 T cell population drives type
1 diabetes. Nature \emph{602}, 156--161. 10.1038/s41586-021-04248-x.}

\href{https://sciwheel.com/work/bibliography/10526050}{40.~  Ellis,
G.I., Sheppard, N.C., and Riley, J.L. (2021). Genetic engineering of T
cells for immunotherapy. Nat. Rev. Genet. \emph{22}, 427--447.
10.1038/s41576-021-00329-9.}

\href{https://sciwheel.com/work/bibliography/18537396}{41.~  Pariset,
M., Hsieh, Y.-P., Bunne, C., Krause, A., and De Bortoli, V. (2023).
Unbalanced Diffusion Schrödinger Bridge. arXiv.
10.48550/arxiv.2306.09099.}

\href{https://sciwheel.com/work/bibliography/770431}{42.~  Kaech, S.M.,
and Wherry, E.J. (2007). Heterogeneity and cell-fate decisions in
effector and memory CD8+ T cell differentiation during viral infection.
Immunity \emph{27}, 393--405. 10.1016/j.immuni.2007.08.007.}

\href{https://sciwheel.com/work/bibliography/15970531}{43.~  Abadie, K.,
Clark, E.C., Valanparambil, R.M., Ukogu, O., Yang, W., Daza, R.M., Ng,
K.K.H., Fathima, J., Wang, A.L., Lee, J., et al. (2024). Reversible,
tunable epigenetic silencing of TCF1 generates flexibility in the T~cell
memory decision. Immunity \emph{57}, 271-286.e13.
10.1016/j.immuni.2023.12.006.}

\href{https://sciwheel.com/work/bibliography/314516}{44.~  Althaus,
C.L., Ganusov, V.V., and De Boer, R.J. (2007). Dynamics of CD8+ T cell
responses during acute and chronic lymphocytic choriomeningitis virus
infection. J. Immunol. \emph{179}, 2944--2951.
10.4049/jimmunol.179.5.2944.}

\href{https://sciwheel.com/work/bibliography/424634}{45.~  De Boer,
R.J., Oprea, M., Antia, R., Murali-Krishna, K., Ahmed, R., and Perelson,
A.S. (2001). Recruitment times, proliferation, and apoptosis rates
during the CD8(+) T-cell response to lymphocytic choriomeningitis virus.
J. Virol. \emph{75}, 10663--10669. 10.1128/JVI.75.22.10663-10669.2001.}

\href{https://sciwheel.com/work/bibliography/8454326}{46.~  Wu, H.,
Kumar, A., Miao, H., Holden-Wiltse, J., Mosmann, T.R., Livingstone,
A.M., Belz, G.T., Perelson, A.S., Zand, M.S., and Topham, D.J. (2011).
Modeling of influenza-specific CD8+ T cells during the primary response
indicates that the spleen is a major source of effectors. J. Immunol.
\emph{187}, 4474--4482. 10.4049/jimmunol.1101443.}

\href{https://sciwheel.com/work/bibliography/18537601}{47.~   McInnes,
L., Healy, J., and Melville, J. (2020). UMAP: Uniform Manifold
Approximation and Projection for Dimension Reduction. arXiv:1802.03426
{[}cs, stat{]}.}

\href{https://sciwheel.com/work/bibliography/147094}{48.~  Wherry, E.J.,
Blattman, J.N., Murali-Krishna, K., van der Most, R., and Ahmed, R.
(2003). Viral persistence alters CD8 T-cell immunodominance and tissue
distribution and results in distinct stages of functional impairment. J.
Virol. \emph{77}, 4911--4927. 10.1128/jvi.77.8.4911-4927.2003.}

\href{https://sciwheel.com/work/bibliography/11131663}{49.~  Milner,
J.J., Toma, C., Quon, S., Omilusik, K., Scharping, N.E., Dey, A.,
Reina-Campos, M., Nguyen, H., Getzler, A.J., Diao, H., et al. (2021).
Bromodomain protein BRD4 directs and sustains CD8 T cell differentiation
during infection. J. Exp. Med. \emph{218}. 10.1084/jem.20202512.}

\href{https://sciwheel.com/work/bibliography/9725107}{50.~  Milner,
J.J., Nguyen, H., Omilusik, K., Reina-Campos, M., Tsai, M., Toma, C.,
Delpoux, A., Boland, B.S., Hedrick, S.M., Chang, J.T., et al. (2020).
Delineation of a molecularly distinct terminally differentiated memory
CD8 T cell population. Proc Natl Acad Sci USA \emph{117}, 25667--25678.
10.1073/pnas.2008571117.}

\href{https://sciwheel.com/work/bibliography/9493658}{51.~  Renkema,
K.R., Huggins, M.A., Borges da Silva, H., Knutson, T.P., Henzler, C.M.,
and Hamilton, S.E. (2020). KLRG1+ Memory CD8 T Cells Combine Properties
of Short-Lived Effectors and Long-Lived Memory. J. Immunol. \emph{205},
1059--1069. 10.4049/jimmunol.1901512.}

\href{https://sciwheel.com/work/bibliography/17508518}{52.~  Fagerberg,
E., Attanasio, J., Dien, C., Singh, J., Kessler, E.A., Abdullah, L.,
Shen, J., Hunt, B.G., Connolly, K.A., De Brouwer, E., et al. (2025).
KLF2 maintains lineage fidelity and suppresses CD8 T cell exhaustion
during acute LCMV infection. Science \emph{387}, eadn2337.
10.1126/science.adn2337.}

\href{https://sciwheel.com/work/bibliography/329098}{53.~  Carlson,
C.M., Endrizzi, B.T., Wu, J., Ding, X., Weinreich, M.A., Walsh, E.R.,
Wani, M.A., Lingrel, J.B., Hogquist, K.A., and Jameson, S.C. (2006).
Kruppel-like factor 2 regulates thymocyte and T-cell migration. Nature
\emph{442}, 299--302. 10.1038/nature04882.}

\href{https://sciwheel.com/work/bibliography/6419040}{54.~  Masopust,
D., and Soerens, A.G. (2019). Tissue-Resident T Cells and Other Resident
Leukocytes. Annu. Rev. Immunol. \emph{37}, 521--546.
10.1146/annurev-immunol-042617-053214.}

\href{https://sciwheel.com/work/bibliography/4577165}{55.~  Milner,
J.J., Toma, C., Yu, B., Zhang, K., Omilusik, K., Phan, A.T., Wang, D.,
Getzler, A.J., Nguyen, T., Crotty, S., et al. (2017). Runx3 programs
CD8+ T cell residency in non-lymphoid tissues and tumours. Nature
\emph{552}, 253--257. 10.1038/nature24993.}

\href{https://sciwheel.com/work/bibliography/1395264}{56.~  Mackay,
L.K., Minnich, M., Kragten, N.A.M., Liao, Y., Nota, B., Seillet, C.,
Zaid, A., Man, K., Preston, S., Freestone, D., et al. (2016). Hobit and
Blimp1 instruct a universal transcriptional program of tissue residency
in lymphocytes. Science \emph{352}, 459--463. 10.1126/science.aad2035.}

\href{https://sciwheel.com/work/bibliography/16725712}{57.~  Shanahan,
S.-L., Kunder, N., Inaku, C., Hagan, N.B., Gibbons, G., Mathey-Andrews,
N., Anandappa, G., Soares, S., Pauken, K.E., Jacks, T., et al. (2024).
Longitudinal intravascular antibody labeling identified regulatory T
cell recruitment as a therapeutic target in a mouse model of lung
cancer. J. Immunol. \emph{213}, 906--918. 10.4049/jimmunol.2400268.}

\href{https://sciwheel.com/work/bibliography/1028483}{58.~  Mackay,
L.K., Rahimpour, A., Ma, J.Z., Collins, N., Stock, A.T., Hafon, M.-L.,
Vega-Ramos, J., Lauzurica, P., Mueller, S.N., Stefanovic, T., et al.
(2013). The developmental pathway for CD103(+)CD8+ tissue-resident
memory T cells of skin. Nat. Immunol. \emph{14}, 1294--1301.
10.1038/ni.2744.}

\href{https://sciwheel.com/work/bibliography/1307038}{59.~  Sheridan,
B.S., Pham, Q.-M., Lee, Y.-T., Cauley, L.S., Puddington, L., and
Lefrançois, L. (2014). Oral infection drives a distinct population of
intestinal resident memory CD8(+) T cells with enhanced protective
function. Immunity \emph{40}, 747--757. 10.1016/j.immuni.2014.03.007.}

\href{https://sciwheel.com/work/bibliography/10765999}{60.~  Walunas,
T.L., Bruce, D.S., Dustin, L., Loh, D.Y., and Bluestone, J.A. (1995).
Ly-6C is a marker of memory CD8+ T cells. J. Immunol. \emph{155},
1873--1883.}

\href{https://sciwheel.com/work/bibliography/9105104}{61.~  Hänninen,
A., Maksimow, M., Alam, C., Morgan, D.J., and Jalkanen, S. (2011). Ly6C
supports preferential homing of central memory CD8+ T cells into lymph
nodes. Eur. J. Immunol. \emph{41}, 634--644. 10.1002/eji.201040760.}

\href{https://sciwheel.com/work/bibliography/4362786}{62.~  Aibar, S.,
González-Blas, C.B., Moerman, T., Huynh-Thu, V.A., Imrichova, H.,
Hulselmans, G., Rambow, F., Marine, J.-C., Geurts, P., Aerts, J., et al.
(2017). SCENIC: single-cell regulatory network inference and clustering.
Nat. Methods \emph{14}, 1083--1086. 10.1038/nmeth.4463.}

\href{https://sciwheel.com/work/bibliography/9127757}{63.~  Van de
Sande, B., Flerin, C., Davie, K., De Waegeneer, M., Hulselmans, G.,
Aibar, S., Seurinck, R., Saelens, W., Cannoodt, R., Rouchon, Q., et al.
(2020). A scalable SCENIC workflow for single-cell gene regulatory
network analysis. Nat. Protoc. \emph{15}, 2247--2276.
10.1038/s41596-020-0336-2.}

\href{https://sciwheel.com/work/bibliography/1306176}{64.~  Mackay,
L.K., Wynne-Jones, E., Freestone, D., Pellicci, D.G., Mielke, L.A.,
Newman, D.M., Braun, A., Masson, F., Kallies, A., Belz, G.T., et al.
(2015). T-box Transcription Factors Combine with the Cytokines TGF-$\beta$ and
IL-15 to Control Tissue-Resident Memory T Cell Fate. Immunity \emph{43},
1101--1111. 10.1016/j.immuni.2015.11.008.}

\href{https://sciwheel.com/work/bibliography/3583222}{65.~  Yu, B.,
Zhang, K., Milner, J.J., Toma, C., Chen, R., Scott-Browne, J.P.,
Pereira, R.M., Crotty, S., Chang, J.T., Pipkin, M.E., et al. (2017).
Epigenetic landscapes reveal transcription factors that regulate CD8+ T
cell differentiation. Nat. Immunol. \emph{18}, 573--582.
10.1038/ni.3706.}

\href{https://sciwheel.com/work/bibliography/1168744}{66.~  Laidlaw,
B.J., Zhang, N., Marshall, H.D., Staron, M.M., Guan, T., Hu, Y., Cauley,
L.S., Craft, J., and Kaech, S.M. (2014). CD4+ T cell help guides
formation of CD103+ lung-resident memory CD8+ T cells during influenza
viral infection. Immunity \emph{41}, 633--645.
10.1016/j.immuni.2014.09.007.}

\href{https://sciwheel.com/work/bibliography/5308382}{67.~  Chou, C.,
Verbaro, D.J., Tonc, E., Holmgren, M., Cella, M., Colonna, M.,
Bhattacharya, D., and Egawa, T. (2016). The Transcription Factor AP4
Mediates Resolution of Chronic Viral Infection through Amplification of
Germinal Center B Cell Responses. Immunity \emph{45}, 570--582.
10.1016/j.immuni.2016.07.023.}

\href{https://sciwheel.com/work/bibliography/13326827}{68.~  White, S.,
Quinn, J., Enzor, J., Staats, J., Mosier, S.M., Almarode, J., Denny,
T.N., Weinhold, K.J., Ferrari, G., and Chan, C. (2021). Flowkit: A
python toolkit for integrated manual and automated cytometry analysis
workflows. Front. Immunol. \emph{12}, 768541.
10.3389/fimmu.2021.768541.}

\href{https://sciwheel.com/work/bibliography/4822624}{69.~  Wolf, F.A.,
Angerer, P., and Theis, F.J. (2018). SCANPY: large-scale single-cell
gene expression data analysis. Genome Biol. \emph{19}, 15.
10.1186/s13059-017-1382-0.}

\href{https://sciwheel.com/work/bibliography/7092677}{70.~  Luecken,
M.D., and Theis, F.J. (2019). Current best practices in single-cell
RNA-seq analysis: a tutorial. Mol. Syst. Biol. \emph{15}, e8746.
10.15252/msb.20188746.}

\href{https://sciwheel.com/work/bibliography/112055}{71.~  Satija, R.,
Farrell, J.A., Gennert, D., Schier, A.F., and Regev, A. (2015). Spatial
reconstruction of single-cell gene expression data. Nat. Biotechnol.
\emph{33}, 495--502. 10.1038/nbt.3192.}

\href{https://sciwheel.com/work/bibliography/7314344}{72.~  Peyré, G.,
and Cuturi, M. (2019). Computational Optimal Transport. FNT in Machine
Learning \emph{11}, 355--206. 10.1561/2200000073.}

\href{https://sciwheel.com/work/bibliography/18541642}{73.~  Cuturi, M.
(2013). Sinkhorn Distances: Lightspeed Computation of Optimal Transport.
In Advances in Neural Information Processing Systems (Curran Associates,
Inc.).}

\href{https://sciwheel.com/work/bibliography/15335444}{74.~  Muzellec,
B., Teleńczuk, M., Cabeli, V., and Andreux, M. (2023). PyDESeq2: a
python package for bulk RNA-seq differential expression analysis.
Bioinformatics \emph{39}. 10.1093/bioinformatics/btad547.}

